\section{Lab Setup and Scenarios}

The testbed consists of two laptops: an \textbf{Acer Aspire A14} and a \textbf{Lenovo} laptop, connected through a \textbf{Vodafone home AP}. Ethernet connections are established via Cat~5 cables; the Acer uses a USB-A to Ethernet adapter.

\textbf{Ethernet Configuration:}
\begin{itemize}
    \item Link capacity: ..Mb/s..
    \item Duplex-mode: ..full/half..
\end{itemize}

\textbf{Wi-Fi Configuration:}
\begin{itemize}
    \item Protocol: ..
    \item Channel: ..
    \item Number of active stations: ..
    \item Nominal max rate: ..
    \item Current data rate: ..
    \item Security protocol: ..
\end{itemize}

\textbf{Network configuration:}
\begin{itemize}
    \item IP address ..
    \item Subnet: 192.168.1.0/24
    \item Gateway: 192.168.1.1 (Vodafone AP)
    \item DNS: ..
\end{itemize}

\textbf{Software:}
Both PCs are running Linux (Ubuntu) and \textit{iperf3} as client and server alternately

We consider the following three scenarios:

\textbf{Scenario 1 – Ethernet to Ethernet:} Both hosts connected via Ethernet (Cat~5) to the AP. This represents the baseline wired configuration.

\textbf{Scenario 2 – WiFi to WiFi:} Both hosts connected wirelessly to the AP. This scenario highlights the overhead introduced by the 802.11 MAC layer.

\textbf{Scenario 3 – WiFi to Ethernet:} One host connected via WiFi, the other via Ethernet. This mixed scenario is representative of typical home network usage.

For each scenario, goodput is measured as a function of the frame size $L \in (0, 1500]$~bytes, using both TCP and UDP. Each experiment is repeated multiple times to obtain statistically significant results. Results are compared against theoretical predictions.
\subsection{UDP configuration and fragmentation}
Because of the datagram-oriented nature of UDP, the payload size chosen in \texttt{iperf3}, using the \texttt{-l} flag, represents the size of each datagram passed to the IP layer. 
However, when a UDP datagram (payload $+$ 8\,B UDP header $+$ 20\,B IP header) exceeds the link MTU (typically 1500\,B for Ethernet), the IP layer must \textbf{fragment} the packet into multiple IP packets. To summarize, this process introduces some risks such as increased loss probability, (if a single fragment is lost the entire datagram is dropped) and reduced efficiency due to extra IP headers.
Finally, to \textbf{avoid fragmentation} in our UDP experiments, we set the payload size to $L = 1472$\,B, so that the total IP packet is exactly $1472 + 8 + 20 = 1500$\,B, matching the Ethernet MTU. The command used for each UDP test was:
\begin{center}
\texttt{iperf3 -c <server> -u -b 2G -l 1472 -t 20}
\end{center}
with the addition of the \texttt{-R} flag for reverse-direction tests (server sends to client). The target bitrate \texttt{-b~2G} was chosen to saturate the link and expose the natural throughput ceiling of each scenario. 
Each test lasted 20 seconds to collect meaningful measurements. Packet captures were collected in parallel using \texttt{tcpdump} for Wireshark analysis.